% 源代码附录所需宏包与命令。
% 将本文件放在 \begin{document} 之前：% 源代码附录所需宏包与命令。
% 将本文件放在 \begin{document} 之前：% 源代码附录所需宏包与命令。
% 将本文件放在 \begin{document} 之前：% 源代码附录所需宏包与命令。
% 将本文件放在 \begin{document} 之前：\input{appendix_code_preamble}

\usepackage{booktabs}
\usepackage{longtable}
\usepackage{xcolor}
\usepackage{fvextra}
\usepackage{hyperref}

\hypersetup{
  colorlinks=true,
  linkcolor=blue!55!black,
  urlcolor=blue!55!black,
  bookmarksopen=true
}

\definecolor{coderule}{RGB}{185,185,185}

% 从 paper/ 目录编译时，项目根目录是“..”。
% 若论文主文件位于其他目录，只需修改这一行。
\newcommand{\ProjectRoot}{..}

% 自动读取一个项目文件并形成独立、可跨页的代码块。
% 参数 1：相对项目根目录的文件名；参数 2：文件作用。
\newcommand{\CodeFile}[2]{%
  \clearpage
  \subsection*{\texttt{\detokenize{#1}}}
  \addcontentsline{toc}{subsection}{\detokenize{#1}}
  \noindent\textbf{文件作用：}#2\par\smallskip
  \VerbatimInput[
    fontsize=\scriptsize,
    numbers=left,
    numbersep=6pt,
    frame=lines,
    framesep=3mm,
    rulecolor=\color{coderule},
    formatcom=\color{black},
    breaklines=true,
    breakanywhere=true,
    breaksymbolleft={},
    tabsize=4
  ]{\ProjectRoot/#1}
}

% 需要手工内嵌某个文件时，将其内容粘贴到 CodeBlock 环境中。
\DefineVerbatimEnvironment{CodeBlock}{Verbatim}{
  fontsize=\scriptsize,
  numbers=left,
  numbersep=6pt,
  frame=lines,
  framesep=3mm,
  rulecolor=\color{coderule},
  breaklines=true,
  breakanywhere=true,
  breaksymbolleft={},
  tabsize=4
}

\setlength{\LTpre}{0.5em}
\setlength{\LTpost}{0.5em}
\setlength{\emergencystretch}{3em}


\usepackage{booktabs}
\usepackage{longtable}
\usepackage{xcolor}
\usepackage{fvextra}
\usepackage{hyperref}

\hypersetup{
  colorlinks=true,
  linkcolor=blue!55!black,
  urlcolor=blue!55!black,
  bookmarksopen=true
}

\definecolor{coderule}{RGB}{185,185,185}

% 从 paper/ 目录编译时，项目根目录是“..”。
% 若论文主文件位于其他目录，只需修改这一行。
\newcommand{\ProjectRoot}{..}

% 自动读取一个项目文件并形成独立、可跨页的代码块。
% 参数 1：相对项目根目录的文件名；参数 2：文件作用。
\newcommand{\CodeFile}[2]{%
  \clearpage
  \subsection*{\texttt{\detokenize{#1}}}
  \addcontentsline{toc}{subsection}{\detokenize{#1}}
  \noindent\textbf{文件作用：}#2\par\smallskip
  \VerbatimInput[
    fontsize=\scriptsize,
    numbers=left,
    numbersep=6pt,
    frame=lines,
    framesep=3mm,
    rulecolor=\color{coderule},
    formatcom=\color{black},
    breaklines=true,
    breakanywhere=true,
    breaksymbolleft={},
    tabsize=4
  ]{\ProjectRoot/#1}
}

% 需要手工内嵌某个文件时，将其内容粘贴到 CodeBlock 环境中。
\DefineVerbatimEnvironment{CodeBlock}{Verbatim}{
  fontsize=\scriptsize,
  numbers=left,
  numbersep=6pt,
  frame=lines,
  framesep=3mm,
  rulecolor=\color{coderule},
  breaklines=true,
  breakanywhere=true,
  breaksymbolleft={},
  tabsize=4
}

\setlength{\LTpre}{0.5em}
\setlength{\LTpost}{0.5em}
\setlength{\emergencystretch}{3em}


\usepackage{booktabs}
\usepackage{longtable}
\usepackage{xcolor}
\usepackage{fvextra}
\usepackage{hyperref}

\hypersetup{
  colorlinks=true,
  linkcolor=blue!55!black,
  urlcolor=blue!55!black,
  bookmarksopen=true
}

\definecolor{coderule}{RGB}{185,185,185}

% 从 paper/ 目录编译时，项目根目录是“..”。
% 若论文主文件位于其他目录，只需修改这一行。
\newcommand{\ProjectRoot}{..}

% 自动读取一个项目文件并形成独立、可跨页的代码块。
% 参数 1：相对项目根目录的文件名；参数 2：文件作用。
\newcommand{\CodeFile}[2]{%
  \clearpage
  \subsection*{\texttt{\detokenize{#1}}}
  \addcontentsline{toc}{subsection}{\detokenize{#1}}
  \noindent\textbf{文件作用：}#2\par\smallskip
  \VerbatimInput[
    fontsize=\scriptsize,
    numbers=left,
    numbersep=6pt,
    frame=lines,
    framesep=3mm,
    rulecolor=\color{coderule},
    formatcom=\color{black},
    breaklines=true,
    breakanywhere=true,
    breaksymbolleft={},
    tabsize=4
  ]{\ProjectRoot/#1}
}

% 需要手工内嵌某个文件时，将其内容粘贴到 CodeBlock 环境中。
\DefineVerbatimEnvironment{CodeBlock}{Verbatim}{
  fontsize=\scriptsize,
  numbers=left,
  numbersep=6pt,
  frame=lines,
  framesep=3mm,
  rulecolor=\color{coderule},
  breaklines=true,
  breakanywhere=true,
  breaksymbolleft={},
  tabsize=4
}

\setlength{\LTpre}{0.5em}
\setlength{\LTpost}{0.5em}
\setlength{\emergencystretch}{3em}


\usepackage{booktabs}
\usepackage{longtable}
\usepackage{xcolor}
\usepackage{fvextra}
\usepackage{hyperref}

\hypersetup{
  colorlinks=true,
  linkcolor=blue!55!black,
  urlcolor=blue!55!black,
  bookmarksopen=true
}

\definecolor{coderule}{RGB}{185,185,185}

% 从 paper/ 目录编译时，项目根目录是“..”。
% 若论文主文件位于其他目录，只需修改这一行。
\newcommand{\ProjectRoot}{..}

% 自动读取一个项目文件并形成独立、可跨页的代码块。
% 参数 1：相对项目根目录的文件名；参数 2：文件作用。
\newcommand{\CodeFile}[2]{%
  \clearpage
  \subsection*{\texttt{\detokenize{#1}}}
  \addcontentsline{toc}{subsection}{\detokenize{#1}}
  \noindent\textbf{文件作用：}#2\par\smallskip
  \VerbatimInput[
    fontsize=\scriptsize,
    numbers=left,
    numbersep=6pt,
    frame=lines,
    framesep=3mm,
    rulecolor=\color{coderule},
    formatcom=\color{black},
    breaklines=true,
    breakanywhere=true,
    breaksymbolleft={},
    tabsize=4
  ]{\ProjectRoot/#1}
}

% 需要手工内嵌某个文件时，将其内容粘贴到 CodeBlock 环境中。
\DefineVerbatimEnvironment{CodeBlock}{Verbatim}{
  fontsize=\scriptsize,
  numbers=left,
  numbersep=6pt,
  frame=lines,
  framesep=3mm,
  rulecolor=\color{coderule},
  breaklines=true,
  breakanywhere=true,
  breaksymbolleft={},
  tabsize=4
}

\setlength{\LTpre}{0.5em}
\setlength{\LTpost}{0.5em}
\setlength{\emergencystretch}{3em}
